% To evaluate the effectiveness of our proposed method, we conducted a series of experiments focusing on OOD rejection and generalization tasks.

% \subsection{Training details}

% We utilize the Dino-v2 ViT backbone trained with registers, using the checkpoints provided by Darcet et al. \cite{darcet2024vision}. Our study employs three sizes of Dino-v2 ViT: Giant, Large, and Base. Keeping the backbone frozen, we then train Linear classifiers (LC) on ImageNet-1k dataset. Our primary focus is on concatenating the CLS with the mean of registers. To demonstrate that registers capture information distinct from patch tokens, we also train LCs on a concatenation of CLS and patch token mean, and on CLS token alone. 


\section{Experiments}
\noindent{\textbf{Datasets}}. We evaluate our proposed approach on the following datasets: \underline{(i) Training/In-Distribution (ID) data}: We train the Linear Classifiers (LC) on ImageNet-1K~\cite{russakovsky2015imagenet} for all experiments. ImageNet-1K is a large-scale imaging benchmark comprising $1.3$ million training images and $50,000$ validation images across $1000$ diverse categories. \underline{(ii) Out-of-Distribution (OOD) generalization benchmarks}: (a) ImageNet-A (In-A)~\cite{hendrycks2021natural} is a dataset comprising $7,500$ adversarially filtered images across a $200$-classes subset of ImageNet-1K’s $1,000$ classes that were found to be challenging to current ImageNet classifiers, (b) ImageNet-R (In-R)~\cite{hendrycks2021many} contains $30,000$ images of different real-world renditions from $200$ classes of ImageNet; (c) ImageNet-S (In-S)~\cite{wang2019learning} contains $50$ black and white sketch images for each class in ImageNet for a total of $50,000$. \underline{(iii) Anomaly rejection benchmarks}: Following standard practice~\cite{energyood}, we consider (a) Texture (DTD), (b) SVHN, (c) Places 365, (d) LSUN, (e) LSUN (Resized) and (f) iSUN datasets. It must be noted that, every image in these datasets is appropriately resized and passed as input into the networks.  

\noindent{\textbf{Setup}}. We utilize the open-source\footnote{https://github.com/facebookresearch/dinov2} Dino-v2 ViT backbones trained with registers~\cite{darcet2024vision} for the experiments involving register tokens, and Dino-v2 ViT backbone trained without registers for other experiments.  We evaluate three variants: ViT Giant, Large, and Base. In all cases, we only train the last linear layer of the network while keeping the backbone frozen. We train the linear layers with the SGD optimizer for $10,000$ iterations, using random-resized-crop data augmentation. The training process was conducted on $16$ NVIDIA V$100$ GPUs, with a batch size of $256$ per GPU and learning rate of $0.01$.

\noindent{\textbf{Metrics}}. We report the Top@1 accuracy to measure OOD generalization performance. For anomaly rejection, we utilize Maximum Softmax Probability (MSP)~\cite{hendrycks2017a} and energy~\cite{energyood} scoring functions and report the FPR@TPR$95$ and AUROC~\cite{energyood} for performance evaluation.  

\noindent{\textbf{Baselines}}. We compare our proposed approach with the following methods to understand the impact of choosing different tokens for training the LC. (i) Concatenation of the [\texttt{CLS}] token and the patch token mean (\texttt{CLS}; $\mu_p$ W/o reg) when the model is trained without registers, as in \cite{oquab2023dinov2}, and (ii) the same concatenation when the model is trained with registers (\texttt{CLS}; $\mu_p$ with reg). To ensure consistent feature scales during concatenation, we use features extracted after the final LayerNorm from the Dinov2 backbones.

\subsection{OOD Generalization}
As shown in Table \ref{tab:results}, our experiments on OOD generalization demonstrate the effectiveness of combining the [\texttt{CLS}] token with the mean of Register tokens (\texttt{CLS}; $\mu_R$). Across all three ViT architectures (Giant, Large and Base), our approach consistently outperforms both the baselines. For ViT-Giant, our method achieves the highest accuracy on all three ImageNet OOD generalization benchmarks. This represents improvements of 2.71\%, 1.46\% and 1.53\% respectively over the best performing baseline of (\texttt{CLS}; $\mu_p$ with reg). Similar trends are observed for ViT-Large and ViT-Base. These results suggest that the combination of by viewing register tokens as auxiliary features and combining it with \texttt{CLS} embeddings yield substantially superior capture more features in terms of generalization and robustness. 
\subsection{Anomaly Rejection}
Our approach demonstrates even significant improvements for the task of anomaly rejection and consistently outperforms other ablations. For instance, with ViT-Base, FPR@TPR95 drops dramatically on SVHN (from 18.22 \% to 12.02\%) and LSUN-R (from 24.85 \% to 17.91 \%). On an average for ViT-Large, our method reduces the mean FPR by 22.18 and 3.22  percentage points compared to the baselines respectively. These improvements hold across all ViT architectures, demonstrating the robustness of our approach. Combining [\texttt{CLS}] and Register tokens enhances the model's ability to detect anomalies by capturing auxiliary information that helps distinguish ID from OOD samples. These findings indicate potential for effective real-world anomaly rejection.

% \subsection{Ablation}

%  We perform an ablation study by concatenating the [\texttt{CLS}] token with all individual register tokens at once (\texttt{CLS}; $r_1$, $r_2$, $r_3$, $r_4$), instead of using the mean of register tokens ($\mu_R$), and train the LC on this representation.
% Table \ref{tab:ablation} shows that while concatenating individual register tokens outperforms the patch token mean, it still falls short of the performance achieved by concatenating the register token mean.

% However, it is possible that concatenating all individual register token may outperform the mean with a more complex fine-tuning approach rather than just linear probing. Nevertheless, exploring this possibility is beyond the scope of our current work.

% \begin{table}[h]
% \caption{Ablation study comparing the performance of different concatenation methods. Concatenating the mean of register tokens (\texttt{CLS}; $\mu_R$) outperforms other approaches.}
% \centering
% \begin{tabular}{|c|ccc|c|}
% \hline
% Method      & \multicolumn{3}{c|}{\begin{tabular}[c]{@{}c@{}}ImageNet OOD\\ Accuracy\end{tabular}} & \begin{tabular}[c]{@{}c@{}}Anomaly Rejection\\ FPR $\downarrow$  / AUROC $\uparrow$ \end{tabular} \\ \hline
%              & In-A                       & In-R                       & In-S                       & Energy; Mean                                                  \\
% \texttt{CLS}; $\mu_P$ & 73.65                      & 76.31                      & 61.86                       & 28.09 / 93.03                                                   \\
% \texttt{CLS}; $\mu_R$ & {\color[HTML]{0f9845}\textbf{78.09}}                      & {\color[HTML]{0f9845}\textbf{80.51}}                     & {\color[HTML]{0f9845}\textbf{64.43}}                      & {\color[HTML]{0f9845}\textbf{20.86 / 94.66}}                                                    \\
% \texttt{CLS}; $r_i$  & 76.57                      & 78.70                      & 62.89                      & 27.63 / 93.63                                                   \\ \hline
% \end{tabular}
% \label{tab:ablation}
% \end{table}